\section{A rank dichotomy for high-girth shores}
\label{sec:four-high-connectivity-entropy-moore}

The following average-growth bound uses a stationary non-backtracking walk.
It avoids an external irregular Moore-bound theorem.
All logarithms and entropies have base two.

\begin{lemma}[An entropy Moore bound for normalized poles]
\label{lem:four-entropy-moore}
Let $N$ be a normalized $k$-pole of rank $h\ge2$, order $n$, and girth $g\ge3$.
Then
\[
 n=2h+k-2,
 \qquad
 g\le4+2\left(1+\frac{k}{3(h-1)}\right)\log n.
\]
\end{lemma}

\begin{proof}
The degree sum gives $2h-2$ cubic vertices and $k$ degree-two vertices.
Hence $n=2h+k-2$ and $m=|E(N)|=3h+k-3$.

Start with a uniformly chosen directed edge.
At each step, choose uniformly among edges that do not reverse the preceding step.
The directed-edge distribution remains uniform.
Indeed, a directed edge leaving $v$ has $\deg(v)-1$ possible predecessors, each assigning it probability $1/(\deg(v)-1)$.
Thus the conditional entropy of each further step is
\[
 \sigma=\frac{1}{2m}\sum_v\deg(v)\log(\deg(v)-1)
       =\frac{3(h-1)}{3h+k-3}>0.
\]

Put $t=\lfloor(g-1)/2\rfloor$.
The random walk of length $t$ has entropy
\[
 H_t=\log(2m)+(t-1)\sigma.
\]
Every such walk is a simple path, since $t<g$.
Two distinct such paths with the same ordered endpoints would give a cycle of length at most $2t<g$.
There are therefore at most $n^2$ possible walks.
Entropy is at most the logarithm of the support size, so
\[
 (t-1)\sigma\le2\log n-\log(2m)\le\log n.
\]
The final inequality uses $2m\ge n$.
Since $t\ge(g-2)/2$, this gives $g\le4+2\log n/\sigma$.
Substitute the formula for $\sigma$.
\end{proof}

\begin{corollary}[Small or polynomial-rank shores]
\label{cor:four-high-girth-shore-dichotomy}
Fix $c,C>0$.
For all sufficiently large $r$, the following holds.
Suppose a normalized $k$-pole has rank $h\ge1$, with
\[
 k\le C\log r,
 \qquad \operatorname{girth}(N)\ge c\log r.
\]
Then either
\[
 h\le1+\frac{4C}{3c}\log(3C\log r),
\]
or
\[
 h\ge r^{c/4}.
\]
The rank threshold depends only on $c,C$.
\end{corollary}

\begin{proof}
The case $h=1$ satisfies the first bound.
Assume $h\ge2$, and write $x=\log r$.
First suppose $h<k$.
Then $n\le3Cx$ and Lemma~\ref{lem:four-entropy-moore} gives
\[
 cx\le4+2\log(3Cx)
       +\frac{2Cx}{3(h-1)}\log(3Cx).
\]
For large $r$, the first two terms sum to at most $cx/2$.
Rearrangement gives the first alternative.

Now suppose $h\ge k$.
Then $n\le3h$ and $k/[3(h-1)]\le2/3$.
Consequently
\[
 cx\le4+\frac{10}{3}\log(3h),
 \qquad
 h\ge\frac{2^{-6/5}}{3}\,r^{3c/10}.
\]
For sufficiently large $r$, the last bound exceeds $r^{c/4}$.
\end{proof}

For a minimum cyclic cut in a cubic graph, every shore is connected and its ports are distinct.
Thus each shore is a normalized pole.
The girth of either shore is at least the ambient girth.
The corollary therefore applies directly when the cut order is $O(\log r)$.

For minimum-cut shores, an elementary argument gives a stronger direct conclusion.
It avoids any external theorem about atoms or cyclic parts.

\begin{corollary}[Every noncycle shore has exponential rank in its girth]
\label{cor:four-minimum-shore-exponential-rank}
Let $N$ be a normalized $k$-pole of rank $h\ge2$ and girth $g\ge\max\{k,3\}$.
Then
\[
 k\le6(h-1),\qquad
 h\ge1+2^{(g-22)/6}.
\]
Consequently, a shore of a minimum cyclic cut in a cubic graph is either a cycle or has rank exponential in the ambient girth.
\end{corollary}

\begin{proof}
Every graph of rank at least two contains two distinct cycles.
If they are edge-disjoint, their union has at least $2g$ edges.
Otherwise their union contains a theta graph.
Its three cycles have total length twice its number of edges, so that number is at least $3g/2$.
In either case, writing $m=3h+k-3$,
\[
 g\le\frac{2m}{3}=2h+\frac{2k}{3}-2.
\]
Since $g\ge k$, it follows that $k\le6(h-1)$.
Therefore
\[
 n=2h+k-2\le8(h-1),\qquad
 \sigma=\frac{3(h-1)}{3(h-1)+k}\ge\frac13.
\]
Lemma~\ref{lem:four-entropy-moore} gives
\[
 g\le4+6\log n\le22+6\log(h-1).
\]
Rearranging proves the exponential bound.

A minimum cyclic-cut shore is a normalized pole, and its girth is at least the ambient girth.
The cut order is at most the ambient girth.
Thus the bound applies whenever its rank is at least two.
A connected normalized pole of rank one has only degree-two vertices, so it is a cycle.
\end{proof}

This is a structural reduction, not a cycle-count bound by itself.
The crossing-count application also needs the growing-terminal forest estimate and its uniform error.

\begin{corollary}[A stronger rank bound and a large port-free ball]
\label{cor:four-near-moore-shore-rank}
Under the hypotheses of Corollary~\ref{cor:four-minimum-shore-exponential-rank}, if $g\ge120$, then
\[
 h\ge2^{g/2-5}.
\]
Moreover, $N$ has a port-free radius-$a$ ball, where
\[
 a=\left\lfloor\frac g2-\log(192g)\right\rfloor.
\]
This ball contains at least $2^{g/2}/(192g)$ vertices.
\end{corollary}

\begin{proof}
The preceding corollary gives $h-1\ge2^{(g-22)/6}\ge g^2$ for $g\ge120$.
The final comparison holds at $120$, and its left-to-right ratio increases thereafter.
If $h\ge2^{g/2}$, the first assertion already holds.
Otherwise
\[
 n=2h+k-2\le2^{g/2+2},\qquad
 \frac{2k\log n}{3(h-1)}
 \le\frac{2g(g/2+2)}{3g^2}<1.
\]
Lemma~\ref{lem:four-entropy-moore} now gives $g<5+2\log n$.
Thus $n>2^{(g-5)/2}$.
For $g\ge120$, one has $g\le2^{(g-5)/2}/2$.
Since $k\le g$,
\[
 h=\frac{n-k+2}{2}>\frac14\,2^{(g-5)/2}>2^{g/2-5}.
\]

Any radius-$a$ ball in a graph of maximum degree three contains at most $3\cdot2^a-2$ vertices.
The union of those balls around all ports therefore has at most
\[
 3k2^a\le3g\frac{2^{g/2}}{192g}=2^{g/2-6}
\]
vertices.
But $n\ge h\ge2^{g/2-5}$, so some vertex lies outside that union.
Its radius-$a$ ball contains no port.
Since $2a<g$, its breadth-first expansion has at least $3\cdot2^a-2\ge2^{a+1}$ vertices.
The definition of $a$ gives $2^{a+1}\ge2^{g/2}/(192g)$, as required.
\end{proof}
